
\section{Conclusion}
In this paper, we presented VideoITG, a novel framework for instruction-aligned frame selection in Video-LLMs. The key to our approach was the \textit{VidThinker} pipeline, which mimics human annotation by generating detailed, instruction-guided clip descriptions, retrieving relevant segments, and performing fine-grained frame selection. Using this pipeline, we constructed the VideoITG-40K dataset with 40K videos and 500K temporal grounding annotations. Based on this resource, we developed plug-and-play VideoITG models that leverage visual-language alignment and reasoning to handle diverse temporal grounding tasks. Experiments showed that VideoITG consistently improves Video-LLMs' performance across multiple video understanding benchmarks, highlighting its effectiveness and potential for advancing instruction-driven video understanding.

% \noindent\textbf{Limitations.} 
% Our current framework consists of two separate modules during inference: VideoITG for frame selection and a standalone Video-LLM for question answering. Although we have carefully designed the frame labeling strategies, the lack of gradient optimization between these two modules during training can lead to suboptimal results. Our VideoITG framework serves as a promising starting point, while in future work we could explore reinforcement learning techniques to bridge these two modules, enabling more efficient and accurate frame selection. 

